\section{Introduction}
\label{introduction}

\subsection{Background and Motivation}
Requirement elicitation determines whether stakeholder goals, domain constraints, and operational assumptions are captured before requirements are formalized.
In intelligent requirements engineering, requirements increasingly become the interface through which stakeholders communicate intent to AI-supported tools~\cite{robinsonRequirementsAreAll2025}.
However, initial natural-language requirements are frequently partial, informal, and context-dependent.

\subsection{Challenges of Requirement Clarification}
LLM-based elicitation can generate fluent requirement statements, but fluency does not guarantee completeness or clarity.
Initial stakeholder inputs may omit actors, data objects, exception cases, quality constraints, or business rules.
Static requirement generation is especially limited because it may rewrite vague input instead of asking targeted follow-up questions.
Prior work on controlled natural language also shows that reducing ambiguity while preserving stakeholder readability remains challenging~\cite{darifControlledNaturalLanguage2026}.

\subsection{Multi-Agent Requirement Clarification}
This paper focuses on multi-agent requirement clarification via missing information detection.
The central idea is to use specialized agents to identify incomplete requirement elements, generate dynamic clarification questions, and keep human stakeholders responsible for approving elicitation decisions.
Knowledge-assisted analysis guides the agents toward functional, constraint, role, data, and exception gaps.

\subsection{Contributions}
The paper is organized around the following contributions:
\begin{itemize}[leftmargin=*]
  \item A research-problem-driven workflow for multi-agent requirement clarification in G1-elicitation.
  \item A missing-information detection scheme covering functional, constraint, role, data, and exception information.
  \item A knowledge-assisted clarification method that uses requirement knowledge, templates, and prompt strategies for dynamic follow-up.
  \item An experimental design for evaluating requirement completeness, missing-information coverage, clarification effectiveness, and human effort.
\end{itemize}
